\section{Proposed Improvements}
\label{sec:method}

We propose three modifications to the \baseline pipeline, each targeting
a specific bottleneck identified through systematic experimentation.
All modifications are implemented as independent toggles and can be
combined freely.

% ------------------------------------------------------------------
\subsection{Diverse ICL Selection via KMeans Clustering}
\label{sec:diverse_icl}

\baseline retrieves the top-$S$ training images by CLIP cosine
similarity.  Because CLIP embeddings cluster tightly by visual
semantics, this often returns near-duplicate images---\eg, ten sunset
photographs for a sunset query---yielding ICL demonstrations that
showcase a single compositional strategy.  The VLM, seeing repetitive
examples, produces a correspondingly narrow distribution of crop
candidates.

We address this by diversifying the ICL set.  We first retrieve $2S$
candidates from the CLIP index, then apply KMeans clustering
($K = S/5$) on their CLIP embeddings to partition them into
semantically distinct groups.  From each cluster, we select
$S/K$ images nearest to the cluster centroid, yielding $S$ total
examples that span a wider range of compositional strategies while
remaining relevant to the query.

This modification does not change the prompt format or VLM
interaction---only the set of images presented as demonstrations.

% ------------------------------------------------------------------
\subsection{Multi-Temperature Candidate Generation}
\label{sec:multitemp}

At the recommended temperature $\tau = 0.05$, the VLM generates $R$
crop candidates that are nearly identical: across our 30-image test
set, the standard deviation of predicted $x_1$ coordinates across the
$R$ candidates for a single image is typically below 5 pixels.  This
means the ``generate $R$, pick the best'' paradigm degenerates into
scoring essentially the same crop $R$ times.

To inject diversity into the candidate pool, we generate crops at two
temperatures: $\tau_{\text{low}} = 0.05$ for precise, high-confidence
proposals, and $\tau_{\text{high}} = 0.8$ for exploratory proposals
that may deviate from the modal prediction.  Each temperature produces
$R$ candidates, and the union of $2R$ candidates is scored and refined
jointly.

This approach is orthogonal to prompt content---it purely expands the
search space over which the scorer operates.

% ------------------------------------------------------------------
\subsection{Expanded Candidate Space}
\label{sec:bigr}

We increase $R$ from 6 to 10, generating 67\% more candidates per VLM
call.  Combined with multi-temperature sampling, this yields up to 20
candidates per refinement iteration, providing a substantially wider
search space for the scorer.  While the original paper~\cite{cropper2025}
reports diminishing returns beyond $R = 5$, this finding was obtained
under single-temperature generation where all candidates are
near-identical.  With diverse generation, additional candidates are
more likely to be genuinely different, improving the returns to
larger~$R$.

% ------------------------------------------------------------------
\subsection{Scorer Debiasing: VILA-Only Scoring}
\label{sec:scorer_debias}

The Area scorer in Eq.~\eqref{eq:scorer} assigns
$\text{Area}(\text{crop}) = A_{\text{crop}} / A_{\text{image}}$, which
monotonically increases with crop size and achieves its maximum
(\ie, 1.0) at the full image.  During iterative refinement, this
creates a positive feedback loop: larger crops receive higher composite
scores, the VLM sees ``large crop scored well'' in the refinement
context, and responds by generating still-larger crops.  We observe
that 90\% of baseline predictions converge to crops covering
$> 95\%$ of the image area.

We resolve this by removing the Area component entirely and scoring
with VILA alone:
\begin{equation}
    \text{Score}_{\text{ours}} = \text{VILA}(\text{crop}).
    \label{eq:scorer_ours}
\end{equation}
VILA evaluates aesthetic quality independent of crop size, breaking the
degenerate feedback loop.

% ------------------------------------------------------------------
\subsection{GAICD Calibration Head}
\label{sec:calhead}

As an additional experiment, we train a lightweight calibration head
that provides a dataset-specific aesthetic signal.  We extract 768-dimensional
features from CLIP ViT-L/14~\cite{clip2021} for each of the ${\sim}288$K
annotated crops in the GAICD training set, and concatenate 6 geometry
features (normalized box coordinates, aspect ratio, area ratio) to form a
774-dimensional input.  A two-layer MLP ($774 \to 256 \to 1$) is trained
to predict MOS via ridge regression.  The calibration head is added as
a third scorer component alongside VILA.

% ------------------------------------------------------------------

% Feedback loop figure
\begin{figure}[t]
\centering
\resizebox{\columnwidth}{!}{%
\begin{tikzpicture}[
    node distance=0.8cm,
    box/.style={draw, rounded corners, minimum height=0.8cm, minimum width=2.0cm,
                align=center, font=\small},
    badbox/.style={box, fill=red!10, draw=red!60},
    goodbox/.style={box, fill=green!10, draw=green!50!black},
    arrow/.style={-{Stealth[length=2.5mm]}, thick},
]
% Bad loop (left)
\node[font=\small\bfseries] (title_bad) {Area Scorer (Baseline)};
\node[badbox, below=0.4cm of title_bad] (vlm1) {VLM generates\\crop candidates};
\node[badbox, below=of vlm1] (area) {Area scorer:\\bigger $=$ better};
\node[badbox, below=of area] (feedback1) {Feedback:\\``big crop scored 0.9''};
\node[badbox, below=of feedback1] (converge) {Converge to\\full image};

\draw[arrow] (vlm1) -- (area);
\draw[arrow] (area) -- (feedback1);
\draw[arrow] (feedback1) -- (converge);
\draw[arrow, dashed, red!60] (feedback1.west) -- ++(-0.6, 0) |- (vlm1.west);

% Good loop (right)
\node[font=\small\bfseries, right=3.5cm of title_bad] (title_good) {VILA-Only (Ours)};
\node[goodbox, below=0.4cm of title_good] (vlm2) {VLM generates\\crop candidates};
\node[goodbox, below=of vlm2] (vila) {VILA scorer:\\aesthetic quality};
\node[goodbox, below=of vila] (feedback2) {Feedback:\\varied scores};
\node[goodbox, below=of feedback2] (diverse) {Meaningful\\refinement};

\draw[arrow] (vlm2) -- (vila);
\draw[arrow] (vila) -- (feedback2);
\draw[arrow] (feedback2) -- (diverse);
\draw[arrow, dashed, green!50!black] (feedback2.east) -- ++(0.6, 0) |- (vlm2.east);

\end{tikzpicture}%
}
\caption{\textbf{Scorer feedback loops.}  \emph{Left:} The Area scorer
creates a degenerate cycle where the refinement loop converges to
full-image crops.  \emph{Right:} VILA-only scoring provides
size-independent aesthetic feedback, enabling meaningful iterative
refinement.}
\label{fig:feedback_loop}
\end{figure}
